\documentclass[fleqn]{article}
\usepackage[utf8]{inputenc}
\usepackage[T2A]{fontenc}
\usepackage[russian]{babel}
\usepackage{graphicx}
\usepackage{amsmath}
\renewcommand{\arraystretch}{1.5}
\renewcommand{\baselinestretch}{1.75}
\usepackage[margin=3cm]{geometry}

\usepackage{amssymb}
\usepackage{caption}
\setlength{\mathindent}{0pt}

\usepackage{titlesec} % уменьшение отступа над главами
\titlespacing*{\section}{0pt}{0pt}{1em}  % уменьшение отступа над главами 



\begin{document}

\title{Теория}
1) Узел - устаревшее название для источника или стока \\
2) Бифуркцация седлового узла - слияние источника и стока и они уничтожаются(?)\\
   1) $ q = p - 1 \quad (\lambda = 1)$ - \textbf{транскритическая бифуркация} (бифуркация седло-узел)\\
   2) $ q = -1 - p \quad (\lambda = -1)$ - \textbf{бифуркация удвоения периода (flip)} (бифуркация удвоения периода)\\
   3) $ q = 1 \quad $ ($\lambda$ - комплексное) - \textbf{бифуркация Неймарка–Сакера} (линии суперкритической бифуркации
Неймарка - Сакерса (бифуркации рождения инвариантной кривой или тора). )\\
J = 1 - консервативная система \\
Для диссипативных систем это не так – они обладают (по крайней мере,
локально) свойством сжатия фазового пространства, в результате чего возмож-
но существование притягивающих объектов – аттракторов, которые отвечают
за установившуюся динамику системы. \\
Аттрактор - «центр притяжения» для движения системы \\
S, J - инварианты, те не меняются при переходе к какой-либо другой косоугольной системе
координат \\
СЗ - мультипликаторы неподвижной точки двумерного отображения\\
$ w = \frac{p}{q} $ - число вращений (не уаерен что это так и что это) \\
\end{document}